%! The Tangent Function — Unit 3, Lesson 3.8 — Guided Notes (Answer Key)
\documentclass[10pt]{article}
\usepackage{precalculus-article}
\usepackage{precalculus-key}   % pulls in -boxes; do NOT also load -boxes
\newcommand{\TallMath}[1]{$\displaystyle #1\rule[-1.4em]{0pt}{3.2em}$}

\begin{document}

\pageheader{Unit 3, Lesson 3.8}{Guided Notes --- Answer Key}
\namedateperiod

\vspace{0.08in}

%----------------------------------------------------------------------
% Objectives
%----------------------------------------------------------------------
\begin{objectivebox}
\textbf{I will be able to\ldots}
\begin{enumerate}[leftmargin=1.5em, itemsep=4pt]
  \item \textbf{LO 3.8.A} --- Construct representations of $f(\theta)=\tan\theta$
    using the unit circle. \ansline{compute $\sin\theta/\cos\theta$ from unit circle coordinates}
  \item \textbf{LO 3.8.B} --- Describe key characteristics of the tangent function
    (period, asymptotes, monotonicity, concavity). \ansline{period $\pi$; asymptotes at $\theta=\pi/2+k\pi$; always increasing; concavity changes at inflection}
  \item \textbf{LO 3.8.C} --- Describe additive and multiplicative transformations
    of the tangent function. \ansline{$g(\theta)=a\tan(b(\theta+c))+d$: dilation $|a|$, period $\pi/|b|$, shift $-c$, translation $d$}
\end{enumerate}
\end{objectivebox}

\vspace{0.06in}

%----------------------------------------------------------------------
% Vocabulary
%----------------------------------------------------------------------
\begin{vocabbox}
\begin{description}[leftmargin=0pt, itemsep=6pt]
  \item[\termblanklong{tangent function}]
    The function $f(\theta) = $ \ans{$\sin\theta/\cos\theta$} that gives the
    \ans{slope} of the terminal ray of angle $\theta$ on the unit circle.

  \item[\termblanklong{vertical asymptote}]
    A vertical line $\theta = a$ where the function values grow without
    \ans{bound}.  For $\tan\theta$, these occur where
    $\cos\theta = $ \ans{$0$}.

  \item[\termblanklong{period (tangent)}]
    The tangent function has period \ans{$\pi$} because the
    slope of the terminal ray repeats every
    \ans{half-revolution ($\pi$ radians)}.
\end{description}
\end{vocabbox}

\vspace{0.06in}

%----------------------------------------------------------------------
% Hook
%----------------------------------------------------------------------
\begin{hookbox}
\textbf{Entry Question:} A lighthouse beam rotates at a constant rate.  The distance
along a straight shoreline to the point illuminated by the beam is a function of the
angle $\theta$ the beam makes with the perpendicular.

\vspace{4pt}
\textbf{Q1.} If the lighthouse is 1 unit from the shore and the angle is $\theta=\pi/4$,
how far along the shore does the beam shine?  (Hint: think about the tangent ratio in a
right triangle.)

\ans{$\tan(\pi/4) = 1$; the beam shines exactly 1 unit along the shore.}

\textbf{Q2.} What happens as $\theta$ approaches $\pi/2$?

\ans{The distance grows without bound; $\tan\theta\to+\infty$ as $\theta\to(\pi/2)^-$.}

\vspace{4pt}
\begin{center}
\begin{tikzpicture}[scale=0.75]
  \draw[->] (-0.3,0) -- (4.2,0) node[right]{\small $\theta$};
  \draw[->] (0,-0.3) -- (0,3.2) node[above]{\small distance};
  \draw[dashed, slate] (3.14,-.2) -- (3.14,3.2)
        node[above,font=\tiny]{$\pi/2$};
  \foreach \x/\l in {1.05/{$\pi/6$}, 2.09/{$\pi/3$}}
        \draw (\x,0.05)--(\x,-0.05) node[below,font=\tiny]{\l};
  % key: draw the tan curve
  \draw[keyred, thick, domain=0.05:2.95, samples=80, smooth]
        plot (\x, {min(3.0,tan(deg(\x*0.5)))});
\end{tikzpicture}
\end{center}

\end{hookbox}

\vspace{0.06in}

%----------------------------------------------------------------------
% LO 3.8.A
%----------------------------------------------------------------------
\begin{notesbox}{LO 3.8.A --- Constructing $\tan\theta$ from the Unit Circle}

\textbf{Definition:}
\[
  \tan\theta = \frac{\ans{$\sin\theta$}}{\ans{$\cos\theta$}}
  \quad\text{provided } \cos\theta \neq \ans{0}
\]

\vspace{6pt}
\textbf{Completed table:}

\vspace{4pt}
\begin{center}
\renewcommand{\arraystretch}{1.55}
\begin{tabular}{|c|c|c|c|}
\hline
$\theta$ & $\sin\theta$ & $\cos\theta$ & $\tan\theta = \sin\theta/\cos\theta$ \\
\hline
$0$           & $0$             & $1$             & \ans{$0$} \\
\hline
$\pi/6$       & $1/2$           & $\sqrt{3}/2$    & \ans{$1/\sqrt{3} = \sqrt{3}/3$} \\
\hline
$\pi/4$       & $\sqrt{2}/2$    & $\sqrt{2}/2$    & \ans{$1$} \\
\hline
$\pi/3$       & $\sqrt{3}/2$    & $1/2$           & \ans{$\sqrt{3}$} \\
\hline
$\pi/2$       & $1$             & $0$             & \ans{undefined} \\
\hline
$2\pi/3$      & $\sqrt{3}/2$    & $-1/2$          & \ans{$-\sqrt{3}$} \\
\hline
$3\pi/4$      & $\sqrt{2}/2$    & $-\sqrt{2}/2$   & \ans{$-1$} \\
\hline
$\pi$         & $0$             & $-1$            & \ans{$0$} \\
\hline
\end{tabular}
\end{center}

\vspace{6pt}
\textbf{Annotated graph:} (asymptotes at $\theta = -\pi/2$ and $\theta = \pi/2$)

\vspace{4pt}
\begin{center}
\begin{tikzpicture}[xscale=1.8, yscale=0.6]
  \draw[->] (-1.8,0) -- (4.0,0) node[right]{\small $\theta$};
  \draw[->] (0,-3.3) -- (0,3.3) node[above]{\small $y$};
  \foreach \x/\l in {-1.571/$-\tfrac{\pi}{2}$, 0/$0$,
                      1.571/$\tfrac{\pi}{2}$, 3.142/$\pi$} {
    \draw (\x,0.08)--(\x,-0.08) node[below,font=\scriptsize]{\l};
  }
  \draw[dashed, slate] (-1.571,-3.2)--(-1.571,3.2);
  \draw[dashed, slate] ( 1.571,-3.2)--( 1.571,3.2);
  \draw[navy, thick, domain=-1.45:-0.0, samples=80, smooth]
        plot (\x, {tan(deg(\x))});
  \draw[navy, thick, domain=0.0:1.45, samples=80, smooth]
        plot (\x, {tan(deg(\x))});
  \draw[navy, thick, domain=1.70:3.10, samples=80, smooth]
        plot (\x, {tan(deg(\x))});
  \foreach \y in {-2,-1,1,2}
    \draw (0.04,\y)--(-0.04,\y) node[left,font=\scriptsize]{$\y$};
  % annotate key points
  \fill[keyred] (0,0)     circle (2pt) node[above right,font=\scriptsize,keyred]{$(0,0)$};
  \fill[keyred] (0.785,1) circle (2pt) node[right,font=\scriptsize,keyred]{$(\pi/4,1)$};
  \fill[keyred] (2.356,-1)circle (2pt) node[right,font=\scriptsize,keyred]{$(3\pi/4,-1)$};
\end{tikzpicture}
\end{center}

\end{notesbox}

\vspace{0.06in}

%----------------------------------------------------------------------
% LO 3.8.B
%----------------------------------------------------------------------
\begin{notesbox}{LO 3.8.B --- Key Characteristics of $y = \tan\theta$}

\renewcommand{\arraystretch}{1.6}
\begin{tabularx}{\linewidth}{@{} p{0.30\linewidth} X @{}}
\hline
\textbf{Characteristic} & \textbf{Description / Value} \\
\hline
Period & \ans{$\pi$} \\
Vertical asymptotes & $\theta = $ \ans{$\pi/2 + k\pi$} for integer $k$ \\
Zeros & $\theta = $ \ans{$k\pi$} for integer $k$ \\
Monotonicity & The function is \ans{increasing} on every interval between consecutive asymptotes \\
Concavity change & Concave \ans{down} for $\theta<0$; concave \ans{up} for $\theta>0$ (within $(-\pi/2, \pi/2)$) \\
\hline
\end{tabularx}

\vspace{8pt}
\textbf{Why period $= \pi$:} After a half-revolution the terminal ray points in the
opposite direction, so both rise and run \ans{change sign},
giving the same \ans{slope/ratio}.

\end{notesbox}

\vspace{0.06in}

%----------------------------------------------------------------------
% LO 3.8.C
%----------------------------------------------------------------------
\begin{notesbox}{LO 3.8.C --- Transformations of $y = \tan\theta$}

For $g(\theta) = a\cdot\tan(b(\theta + c)) + d$:

\vspace{4pt}
\renewcommand{\arraystretch}{1.6}
\begin{tabularx}{\linewidth}{@{} p{0.12\linewidth} p{0.32\linewidth} X @{}}
\hline
\textbf{Param.} & \textbf{Effect} & \textbf{Example} \\
\hline
$a$ & Vertical dilation by \ans{$|a|$} & $2\tan\theta$: outputs $\times 2$ \\
$b$ & Period becomes \ans{$\pi/|b|$} & $\tan(2\theta)$: period $= \pi/2$ \\
$c$ & Phase shift of \ans{$-c$} & $\tan(\theta-\pi/4)$: shift right $\pi/4$ \\
$d$ & Vertical translation by \ans{$d$} & $\tan\theta + 3$: shift up $3$ \\
\hline
\end{tabularx}

\vspace{8pt}
\textbf{Finding new asymptotes:} Asymptotes occur when the \emph{inside argument}
equals $\pm\pi/2 + k\pi$.

\vspace{4pt}
\textbf{Example:} Find the asymptotes of $g(\theta) = \tan(2\theta)$.

Set $2\theta = \dfrac{\pi}{2} + k\pi$ \quad$\Rightarrow$\quad $\theta = $ \ans{$\dfrac{\pi}{4} + \dfrac{k\pi}{2}$}

\end{notesbox}

\vspace{0.06in}

%----------------------------------------------------------------------
% Practice
%----------------------------------------------------------------------
\begin{practicebox}

\textbf{Guided Practice.}\quad For each function below, state the period and
the location of the vertical asymptote closest to $\theta = 0$ (with $\theta > 0$).

\vspace{4pt}
\renewcommand{\arraystretch}{1.6}
\begin{tabularx}{\linewidth}{@{} p{0.40\linewidth} X X @{}}
\hline
\textbf{Function} & \textbf{Period} & \textbf{Asymptote ($\theta > 0$)} \\
\hline
$g(\theta) = \tan(3\theta)$ & \ans{$\pi/3$} & \ans{$\theta = \pi/6$} \\
$g(\theta) = \tan\!\left(\tfrac{1}{2}\theta\right)$ & \ans{$2\pi$} & \ans{$\theta = \pi$} \\
$g(\theta) = 4\tan(\theta) + 1$ & \ans{$\pi$} & \ans{$\theta = \pi/2$} \\
\hline
\end{tabularx}

\begin{teachernote}
The $a$ parameter does not change the period or asymptote locations; only $b$ does.
Emphasize to students that the vertical dilation only affects output values.
\end{teachernote}

\end{practicebox}

\end{document}
